\documentclass[a4paper]{article}
\usepackage{amsfonts}
\usepackage{amsmath}
\usepackage{amssymb}
\usepackage{graphicx}     

\begin{document}
  
\noindent \large Auteur: Dara van Engelen \\
\noindent \large Studierichting: Informatica\\
\noindent \large Oefeningenreeks 3F nr. 1.1\\

\medskip

\normalsize

$f(x)=\dfrac{x-1}{x^2+x-6}=\dfrac{x-1}{(x-2)(x+3)}$\\

Verticale asymptoten:\\

$\lim\limits_{x \rightarrow 2+}f(x)=+\infty$ en $\lim\limits_{x \rightarrow 2-}f(x)=-\infty\Rightarrow x = 2$ is een V.A.\\

$\lim\limits_{x \rightarrow -3+}f(x)=+\infty $ en $\lim\limits_{x \rightarrow -3-}f(x)=-\infty \Rightarrow x= -3$ is ook een V.A.\\

Horizontale asymptoten:\\

$\lim\limits_{x \rightarrow \infty}\dfrac{x-1}{x^2+x-6}=0\Rightarrow y = 0$ is een horizontale asymptoot.\\

Geen schuine asymptoot want we hebben al een verticale en horizontale.\\

\begin{tabular}{c|ccccccccc}
$x$			& & $-3$ & & $2$ &\\ \hline
$f(x)-0$ 		& $-$ &   $|$   & $+$ & $|$ & $+$\\ 
			& onder H.A. & & boven H.A. & & boven H.A.\\
\end{tabular}

\begin{thebibliography}{999}
%\bibitem{a} Referentie
\end{thebibliography}
\end{document} 